\documentclass{article}

\usepackage{amsmath, amsthm, amssymb, amsfonts}
\usepackage{thmtools}
\usepackage{graphicx}
\usepackage{setspace}
\usepackage{geometry}
\usepackage{float}
\usepackage{hyperref}
\usepackage[utf8]{inputenc}
\usepackage[english]{babel}
\usepackage{framed}
\usepackage[dvipsnames]{xcolor}
\usepackage{tcolorbox}
\usepackage{physics}

\DeclareMathOperator{\grd}{grad}
\DeclareMathOperator{\crl}{curl}
\DeclareMathOperator{\dvr}{div}
\colorlet{LightGray}{White!90!Periwinkle}
\colorlet{LightOrange}{Orange!15}
\colorlet{LightGreen}{Green!15}

\newcommand{\HRule}[1]{\rule{\linewidth}{#1}}

\declaretheoremstyle[name=Theorem,]{thmsty}
\declaretheorem[style=thmsty,numberwithin=section]{theorem}
\tcolorboxenvironment{theorem}{colback=LightGray}

\declaretheoremstyle[name=Proposition,]{prosty}
\declaretheorem[style=prosty,numberlike=theorem]{proposition}
\tcolorboxenvironment{proposition}{colback=LightOrange}

\declaretheoremstyle[name=Principle,]{prcpsty}
\declaretheorem[style=prcpsty,numberlike=theorem]{principle}
\tcolorboxenvironment{principle}{colback=LightGreen}

\setstretch{1.2}
\geometry{
    textheight=9in,
    textwidth=6in,
    top=0.5in,
    headheight=12pt,
    headsep=25pt,
    footskip=30pt
}

% ------------------------------------------------------------------------------

\begin{document}


% ------------------------'

\textbf{HW 9.}

\textbf{Exercise 1.} Fix some $\alpha$ with $0 < \alpha < \pi/2$ and some $h>0$. Let 
\[X(\theta,t) = (t\sin \alpha \cos \theta, t \sin \alpha \sin \theta, t \cos \alpha), \]
where $\theta \in [0,2\pi]$ and $t \in [0,h\sec \alpha]$. This parametrizes a cone of height $h$. Compute the
tangent vectors $\partial X / \partial \theta$ and $\partial X/ \partial t$ as well as their cross product.
Find the equation for the surface in rectangular/Cartesian coordinates. 

\textbf{Exercise 2.} Fix $h, a >0$ Let 
\[X(t,\theta) = (a t \cos \theta, a t \sin \theta, t^2)\]
where $\theta \in [0,2\pi]$ and $t \in [0,h]$. This parametrizes a paraboloid. Compute the
tangent vectors $\partial X / \partial \theta$ and $\partial X/ \partial t$ as well as their cross product.
Find the equation for the surface in rectangular/Cartesian coordinates. What surface do we obtain if we replace
$t^2$ with $t$ in the $z$ coordinate?

\textbf{Exercise 3.} Compute the surface area of the torus parametrized by 
\[X(\theta, \varphi) = \left((a+b\cos \varphi)\cos\theta, (a+b\cos\varphi)\sin\theta, b \sin \varphi\right)\]
where $0 \leq \varphi < 2\pi$ and $0 \leq \theta < 2\pi$. (Hint: the cross product computation is a little messy,
but using the identity $\sin^2 x + \cos^2 x = 1$ a couple of times, you can simplify the expression for $\norm{\frac{\partial X}{\partial \varphi} \times \frac{\partial X}{\partial \theta}}$ 
considerably).

\textbf{Exercise 4.} Let $\varphi: \mathbb{R}^3 \to \mathbb{R}$ be a smooth function. Show that
\[ \crl (\grd \varphi) = 0. \]

\textbf{Exercise 5.} Let $F: \mathbb{R}^3 \to \mathbb{R}^3$ be a smooth vector field. Show that $\dvr \crl F = 0$.

\textbf{Exercise 6.} Compute the integral 
\[\int_S \crl F \cdot \vec{n} d\sigma\]
where $F = (-y,x^2,z^3)$, and $S$ is the surface
\[x^2 + y^2 + z^2 = 1,\ -\frac{1}{2} \leq z \leq 1.\]
[Hint: by Stokes' theorem, the integral above of any two surfaces that share the same boundary will have the same
value, so one can deform this surface (while holding the boundary fixed) to a much simpler one and compute the integral of $\crl F$ on the deformed surface instead.]

\textbf{Exercise 7.} Let $C$ be a closed curve which is the 
boundary of a surface $S$. Let $f$ and $g$ be scalar-valued functions. Show that
\[\int_C (f \grd g) \cdot dC = \iint_S [(\grd f) \times (\grd g)]\cdot \mathbf{n} d \sigma.\]
Show consequently that 
\[\int_C (f \grd g + g \grd f)\cdot dC = 0.\]


\end{document}